%%%%%%%%%%%%%%%%%%%%%%%%%%%%%%%%%%%%%%%%%%%%%%%%%%%%%%%%%%%%%%%%%%%%%%%%%%%%%%%
%% Dados do trabalho
%%%%%%%%%%%%%%%%%%%%%%%%%%%%%%%%%%%%%%%%%%%%%%%%%%%%%%%%%%%%%%%%%%%%%%%%%%%%%%%

%% Título definido pelo autor em 08/11/2026, entre as alternativas de uma pesquisa
%% sobre títulos (NBR 14724:2024, 4.1.1 c e d; guia de redação da USC; estudos
%% sobre tamanho e tipo de título). O título nomeia o objeto e o método; o
%% subtítulo, depois de dois-pontos, nomeia a operação feita (avaliação) e o
%% objeto avaliado (o sistema). Sem resultado, sem sigla e sem nome de projeto.
%% A quebra depois dos dois-pontos separa título e subtítulo nas três folhas; nos
%% metadados do PDF ela vira espaço.
\titulo{Detecção de falhas em máquinas industriais por afastamento da operação normal:\texorpdfstring{\\}{ }avaliação de um sistema de baixo custo}
\autor{Lauan Matheus da Rocha Alves}
\local{Teresina -- PI}
\data{2026}
\orientador[Orientadora:]{Lia Mariana Nascimento Brito}
%\coorientador{Nome do coorientador}
\preambulo{Trabalho de Conclusão de Curso apresentado como requisito parcial
para obtenção do título de Bacharel em Engenharia de Software do Instituto de
Ensino Superior -- iCEV.}

%% Metadados do PDF: precisam vir depois de \titulo e \autor.
\makeatletter
\hypersetup{pdftitle = {\@title}, pdfauthor = {\@author}}
\makeatother

%% Usados na capa do modelo iCEV (preambulo/estilo-icev.tex).
\newcommand{\cursoicev}{Curso de Bacharelado em Engenharia de Software}
\newcommand{\instituicaoicev}{Instituto de Ensino Superior -- iCEV}
